\section{Check the reconstructed integer before narrowing it}
A canonical mapped value does not by itself establish that the next reading fits.
Apply the inverse mapping in a doubleword, add the previous reading in that same width,
and compare the exact result with \(-2147483648\) and 2147483647. Only then store the
low word, whose bits now represent a proven signed 32-bit result.

For bytes \texttt{fe ff ff ff 0f 02}, the first five bytes reconstruct the maximum
reading. The next mapped value two denotes a delta of one. Its exact new reading is
2147483648, so return range status three with five committed input bytes and one
output reading. A word addition would wrap it to the minimum reading and accept data
that the format rejects.

One implementation of the inverse uses an unsigned right shift by one and an XOR with
the negative of the low bit. Another branches on parity and computes
\(-\lfloor u/2\rfloor-1\) for odd values. Both must retain the full unsigned value
until the sign decision is complete. A signed shift of a truncated word can lose the
fifth byte's contribution before any range check occurs.

\begin{MARSExercise}{Separate two error categories}
Compare \texttt{80 80 80 80 20} and \texttt{ff ff ff ff 1f}. Both fail to produce a
reading from predecessor zero. State why the first reports malformed status two and
the second reports range status three. Explain which check takes priority in each case.
\end{MARSExercise}
